% ch08-payments-jobs — الحجز والدفع (Money Flow) والمهامّ الخلفيّة
\section{الحجز والدفع (Money Flow) والمهامّ الخلفيّة}

في هذا الفصل نمشي مع «الفلوس» من أوّل ما الطالب يحجز جلسة مع مستشار، لحدّ ما
تصل صافي أرباحها لجيب المستشار. الفكرة الأساسية بسيطة جدًّا في جملة واحدة:
\textbf{فلوس الطالب تتحجز ولا تتاخد}، ولا تُخصَم فعليًّا إلا لمّا المستشار يقبل
الجلسة. ده نموذج شبيه بالـ \en{escrow} (الوسيط الأمين) مبنيّ على \en{Stripe}.
وفي آخر الفصل هنشوف «المهامّ الخلفيّة» (\en{background jobs}) اللي بتشتغل لوحدها
في الخلفية: تذكيرات، صرف أرباح، تنظيف، والتزامات \en{GDPR}.

كل الكلام اللي جاي مقتبس من الكود الحقيقي للنظام في \code{server/src/...}؛
هنعرض مقتطفات قصيرة حقيقية ونشرحها سطرًا بسطر.

\subsection{لماذا نموذج authorize-then-capture؟}

تخيّل إنّ الطالب دفع كامل المبلغ لحظة الحجز، وبعدين المستشار رفض أو مردّش. دلوقتي
إحنا ماسكين فلوس الطالب ولازم نعمل عملية استرداد كاملة. الأفضل من كده إنّنا
\textbf{نحجز} المبلغ على البطاقة (نأخذ «تفويضًا» فقط) من غير ما نخصمه، ونستنّى.

\begin{idea}[الفكرة: التفويض ثم القبض (authorize then capture)]
عملية الدفع على البطاقة بتتقسّم لخطوتين منفصلتين:
\begin{itemize}
\item \textbf{التفويض} (\en{authorize}): البنك بيحجز المبلغ على بطاقة الطالب
  ويضمن إنّه موجود، لكن الفلوس \emph{لسه ما اتحوّلتش} لحسابنا.
\item \textbf{القبض} (\en{capture}): دي الخطوة اللي بتحوّل الفلوس فعليًّا. إحنا
  بنأجّلها لحدّ ما المستشار يقبل الجلسة.
\end{itemize}
ما بين الخطوتين، فلوس الطالب «محجوزة» في مكان آمن: لا هو خسرها، ولا إحنا أخدناها.
\end{idea}

المنفذ (\en{port}) اللي بيمثّل كل تعاملاتنا مع \en{Stripe} اسمه \code{IStripeService}،
وله محوّلان (\en{adapters}): \code{StripeService} (الحقيقي، بيتكلّم مع \en{Stripe}
فعلًا) و\code{StubStripeService} (بديل وهميّ للتطوير والاختبار، بيردّ نتائج جاهزة
من غير أي اتّصال بالإنترنت). زيّ ما اتكلّمنا في مبدأ الـ \en{ports} قبل كده، ده
بيخلّي منطق الحجز نفسه ما يعرفش حاجة عن \en{Stripe}، ونقدر نبدّل المزوّد أو نختبر
من غير ما نلمس قواعد العمل.

\begin{why}
ليه التفويض-ثم-القبض ومش الخصم المباشر؟ لأنّه بيحمي الطرفين. الطالب مش هيتخصم منه
قرش إلا لو الخدمة اتأكّدت فعلًا (المستشار قبل)، والمنصّة مش محتاجة تعمل عمليات
استرداد كتير ومعقّدة لو الحجز اتلغى قبل ما يبدأ. ده كمان بيقلّل نزاعات البطاقات
(\en{chargebacks}) لأنّ المبلغ ما بيتقبضش غير مقابل خدمة حقيقية.
\end{why}

\subsection{دورة الحياة خطوةً بخطوة}

خلّينا نتفرّج على رحلة الدفعة الواحدة من أوّلها لآخرها. الجدول ده بيلخّص المحطّات
الخمس، وبعده هنفصّل كل واحدة بالكود:

\begin{center}
\begin{tabular}{L{3.6cm} L{4.2cm} L{5.2cm}}
\toprule
\textbf{\color{spPrimary}المحطّة} & \textbf{\color{spPrimary}الفعل التقني} & \textbf{\color{spPrimary}مصير الفلوس} \\
\midrule
حجز + تفويض & \en{PaymentIntent} بـ \en{manual capture} & محجوزة على البطاقة (لا مخصومة) \\
المستشار يقبل & \code{CapturePaymentIntentAsync} & تُقبَض فعليًّا وتدخل النظام \\
رفض / انتهاء المهلة & \code{CancelPaymentIntentAsync} & يُفَكّ الحجز، الطالب ما اتخصمش \\
\en{no-show} / إلغاء & \code{RefundPaymentAsync} + \code{RefundCalculatorService} & استرداد كامل أو جزئيّ حسب السياسة \\
صرف للمستشار & \code{StripePayoutJob} + \en{Stripe Connect} & صافي الأرباح يوصل للمستشار \\
\bottomrule
\end{tabular}
\end{center}

\subsubsection{المحطّة الأولى: الحجز والتفويض}

هنشوف إزاي بننشئ الـ \en{PaymentIntent} بنمط \en{manual capture} لحظة الحجز.
ده مقتطف من إنشاء نيّة الدفع في \code{StripeService}:

\begin{english}
\begin{lstlisting}[language={[Sharp]C}]
var createOptions = new PaymentIntentCreateOptions
{
    Amount = amountCents,           // integer cents — no floating point
    Currency = currency,            // "usd"
    CaptureMethod = captureMethod,  // "manual" -> authorize, don't charge
    Metadata = metadata?.ToDictionary(k => k.Key, v => v.Value),
    PaymentMethodTypes = ["card"],  // card only
};

var requestOptions = new RequestOptions
{
    ApiKey = _opts.SecretKey,
    IdempotencyKey = idempotencyKey, // retry-safe: never double-charge
};

var intent = await service.CreateAsync(createOptions, requestOptions, ct);
\end{lstlisting}
\end{english}

نفصّص أهمّ الأسطر:
\begin{itemize}
\item \code{Amount = amountCents}: المبلغ دايمًا بالسنتات كعدد صحيح
  (\en{integer cents}). يعني 10.50 دولار بتتخزّن 1050. ليه؟ عشان نتجنّب أخطاء
  الكسور العشرية في أرقام الفاصلة العائمة (\en{floating point}). قرش واحد
  ضايع في حسابات الفلوس مشكلة كبيرة.
\item \code{CaptureMethod = "manual"}: ده قلب الموضوع كلّه. «\en{manual}» معناها
  «فوّض دلوقتي، واقبض لمّا أنا أقولك». من غير السطر ده، \en{Stripe} كان هيخصم
  فورًا.
\item \code{PaymentMethodTypes} بقيمة \code{card}: نقصر الدفع على الكروت بس. من
  غير كده \en{Stripe} بيورّي كل المحافظ المفعّلة في لوحة التحكّم وده كان بيلخبط
  المستخدمين.
\item \code{IdempotencyKey}: مفتاح «عدم التكرار». لو الطلب اتبعت مرّتين (بسبب
  ضغطة زرّ مكرّرة أو إعادة محاولة الشبكة)، \en{Stripe} بيتعرّف على المفتاح ده
  ويرجّع نفس النتيجة \emph{من غير ما يعمل حجزًا تاني}. الطالب ما بيتخصمش مرّتين.
\end{itemize}

الطريف إنّ صفّ \code{Payment} في قاعدة بياناتنا بيبدأ بحالة \code{Pending} (مش
\code{Held})، لأنّ نيّة الدفع اليدوية «مش محجوزة فعلًا» إلا بعد ما الطالب يفوّض
البطاقة في المتصفّح. ساعتها \en{Stripe} بيبعت لنا \en{webhook} يقول «تمّ التفويض»،
فنقلب الحالة لـ \code{Held}. (هنرجع للـ \en{webhooks} بعد شوية.)

\subsubsection{المحطّة الثانية: القبول يعني القبض}

لمّا المستشار يقبل الحجز، بنستدعي \code{CapturePaymentIntentAsync} عشان نقبض
المبلغ المحجوز. لو النيّة مش في حالة «جاهزة للقبض» (يعني الطالب ما فوّضش أصلًا)،
بنترجم الخطأ لرسالة واضحة بدل خطأ 500 غامض:

\begin{english}
\begin{lstlisting}[language={[Sharp]C}]
catch (StripeException ex)
{
    // A capture failure is almost always a payment-state problem — the
    // student never authorized the card, so the intent is not in
    // requires_capture. Surface it as a clean 422, not an opaque 500.
    throw new BookingDomainException(
        "The session fee could not be captured. The student may not have "
        + "completed the card authorization for this booking yet — ask them "
        + "to finish payment, then try accepting again.",
        ex);
}
\end{lstlisting}
\end{english}

هنا الفكرة: لو القبض فشل، غالبًا السبب إنّ الطالب ما كمّلش التفويض. فبدل ما
نرمي خطأ سيرفر مبهم، بنرمي استثناء عمل (\en{domain exception}) بيتحوّل لردّ 422
برسالة مفهومة للمستشار: «كمّل الطالب الدفع الأوّل، وبعدين اقبل تاني». تجربة أفضل
بكتير من شاشة خطأ حمرا.

\subsubsection{المحطّة الثالثة والرابعة: الإفراج والاسترداد}

لو المستشار رفض الحجز أو انتهت مهلة ردّه، بنفكّ الحجز عبر
\code{CancelPaymentIntentAsync}، والطالب ما بيتخصمش أصلًا. أما لو حصل إلغاء بعد
القبول، أو حصل \en{no-show} (حد ما حضرش)، فالموضوع محتاج حساب: كام نرجّع للطالب؟
ده شغل \code{RefundCalculatorService}. خلّينا نشوف منطق الطالب اللي بيلغي:

\begin{english}
\begin{lstlisting}[language={[Sharp]C}]
var moreThan24HoursBefore = scheduledStartAt.ToUniversalTime() > nowUtc.AddHours(24);

if (moreThan24HoursBefore)
{
    return new RefundCalculationResult
    {
        RefundPercentage = 100,
        RefundAmountCents = amountCents,
        CancellationReason = CancellationReason.StudentCancelledMoreThan24HoursBefore
    };
}

return new RefundCalculationResult
{
    RefundPercentage = 50,
    RefundAmountCents = amountCents / 2,
    CancellationReason = CancellationReason.StudentCancelledLessThan24HoursBefore
};
\end{lstlisting}
\end{english}

القاعدة زيّ ما بتشوف مباشرة وعادلة:
\begin{itemize}
\item لو الطالب لغى قبل الموعد بأكتر من 24 ساعة: استرداد كامل 100\%.
\item لو لغى في آخر 24 ساعة: استرداد 50\% بس (لأنّ المستشار حجز وقته).
\item لو \emph{المستشار} هو اللي لغى: الطالب بياخد 100\% دايمًا (ولو كان في آخر
  24 ساعة، بيتّطبق كمان خصم على تقييم المستشار من قِبَل المُعالِج).
\end{itemize}

\begin{keyp}[نقطة مهمة: أساس الاسترداد هو المبلغ المخزَّن]
لاحظ إنّ الحساب بيستخدم \code{amountCents} المخزَّن على صفّ الـ \code{Payment}،
مش سعر الحجز الحاليّ. ليه؟ لأنّ السعر ممكن يتغيّر بعد كده (خصم، عملة، إعادة
تسعير)، فاللقطة (\en{snapshot}) المخزَّنة وقت الدفع هي المرجع الموثوق. ده مبدأ
مهمّ: في المحاسبة، بنعتمد على القيمة وقت المعاملة مش القيمة الحالية.
\end{keyp}

بعد ما الحاسبة تقرّر المبلغ، بيتنفّذ الاسترداد فعليًّا عبر \code{RefundPaymentAsync}
على \en{Stripe}. ولو الاسترداد بسبب \en{no-show}، فالمهمّة الخلفيّة
\code{MeetingNoShowSweepJob} هي اللي بترصد الموقف وتشغّل المسار ده تلقائيًّا.

\subsection{تقسيم الأرباح}

كل صفّ \code{Payment} بيخزّن ثلاثة أرقام رئيسية بالسنتات:

\begin{center}
\begin{tabular}{L{4.4cm} L{8.4cm}}
\toprule
\textbf{\color{spPrimary}الحقل} & \textbf{\color{spPrimary}معناه} \\
\midrule
\code{AmountCents} & إجمالي المبلغ اللي دفعه الطالب (الجُملة). \\
\code{ProfitShareAmountCents} & حصّة المنصّة (الرسوم + نسبة الربح). \\
\code{PayeeAmountCents} & صافي نصيب المستشار بعد خصم حصّة المنصّة. \\
\bottomrule
\end{tabular}
\end{center}

النِّسَب نفسها مش مكتوبة في الكود؛ هي بتيجي من قاعدة ماليّة قابلة للتهيئة اسمها
\code{FinancialConfigRule} (ومعاها \code{ProfitShareConfig}). الحساب نفسه دالة
رياضية نقيّة (\en{pure function}) بتشتغل بالسنتات الصحيحة عشان المجموع يرجع
مظبوط بالظبط:

\begin{english}
\begin{lstlisting}[language={[Sharp]C}]
var feeCents = feeKind switch
{
    FeeKind.Percentage => (long)Math.Round(
        grossAmountCents * (feePercentage ?? 0m), MidpointRounding.AwayFromZero),
    FeeKind.FixedAmount => feeAmountCents ?? 0L,
    _ => 0L,
};

var profitShareCents = (long)Math.Round(
    grossAmountCents * profitSharePercentage, MidpointRounding.AwayFromZero);

var platformTotalCents = feeCents + profitShareCents;
// Payee takes the remainder so the split always re-sums to gross exactly.
var payeeNetCents = grossAmountCents - platformTotalCents;
\end{lstlisting}
\end{english}

السطر الأخير هو الحيلة الذكية: نصيب المستشار = الإجمالي ناقص حصّة المنصّة.
بكده مهما حصل تقريب في الكسور، \textbf{مجموع الحصّتين بيساوي الإجمالي بالظبط}
ولا يضيع قرش. الرسوم ممكن تكون نسبة مئوية أو مبلغًا ثابتًا (حسب \code{FeeKind})،
لكن نسبة الربح دايمًا نسبة مئوية.

\begin{note}
النِّسَب بتتحدّد عند القبول، مش عند الحجز. عند الحجز بيتخزّن
\code{ProfitShareAmountCents} بصفر و\code{PayeeAmountCents} بكامل المبلغ (كأنّ
كل المبلغ للمستشار مؤقّتًا)، وبس لمّا المستشار يقبل والقبض ينجح، بيتعاد حساب
الحصّتين من القاعدة الماليّة السارية ساعتها. كده لو الأدمن غيّر النِّسَب، الحجوزات
القديمة المدفوعة ما تتأثّرش.
\end{note}

\subsubsection{قاعدة فعّالة واحدة لكل نوع}

طب لو في أكتر من قاعدة ماليّة لنفس نوع الدفع، نستخدم أنهي واحدة؟ الحلّ أنيق:
\textbf{فهرس فريد مُصفّى} (\en{filtered unique index}) في قاعدة البيانات بيضمن إنّه
لا يمكن وجود أكتر من قاعدة \emph{فعّالة} واحدة لكل نوع:

\begin{english}
\begin{lstlisting}[language={[Sharp]C}]
// FR-170: at most one Active rule per payment type.
b.HasIndex(r => r.PaymentType)
    .IsUnique()
    .HasFilter("[Status] = 'Active'")
    .HasDatabaseName("UX_FinancialConfigRule_ActivePerType");
\end{lstlisting}
\end{english}

الشرط \code{HasFilter} بقيمة \code{[Status] = 'Active'} معناه إنّ قيد التفرّد
بيتطبّق بس على الصفوف الفعّالة. يعني ممكن يكون عندك عشر قواعد قديمة بحالة غير
فعّالة، لكن \emph{واحدة بس} فعّالة لكل نوع. لو حد حاول يفعّل تانية، قاعدة البيانات
نفسها بترفض. القاعدة دي محفورة في الداتابيز مش في كود التطبيق، فمستحيل تتنسى.

\subsection{الصرف للاستشاري عبر Stripe Connect}

لحدّ دلوقتي، الفلوس المقبوضة موجودة في حساب المنصّة على \en{Stripe}. إزاي توصل
للمستشار؟ عن طريق \en{Stripe Connect}: كل مستشار بيعمل حساب \en{Express} خاص بيه
(عبر \code{CreateConnectAccountAsync} + رابط تسجيل)، وبعدين مهمّة خلفيّة اسمها
\code{StripePayoutJob} بتشتغل كل ليلة وتجمّع كل الدفعات المقبوضة اللي لسه ما
اتصرفتش لكل مستشار، وتبعتله دفعة واحدة (\en{payout}) بمجموع \code{PayeeAmountCents}
بتاعته.

\begin{idea}[الفكرة: المطالبة المُسبقة قبل استدعاء Stripe]
مهمّة الصرف بتعمل حاجة ذكية: قبل ما تنادي \en{Stripe} خالص، بتسجّل صفّ الـ
\code{Payout} في قاعدة البيانات وتربط الدفعات بيه (\code{PayoutId}). ليه الأول؟
عشان لو المهمّة اتشغّلت مرّتين بالغلط، النسخة التانية تلاقي الدفعات «متطالَب بيها»
فما تصرفهاش تاني. ولو استدعاء \en{Stripe} فشل بعد كده، بيتشال الربط
(بترجع \code{PayoutId} فاضية) عشان الليلة الجاية تحاول من جديد. وكمان الصرف نفسه
بيتبعت بمفتاح عدم تكرار مشتقّ من معرّف الـ \code{Payout}، فمفيش صرف مزدوج.
\end{idea}

\subsection{مطابقة الـ webhook والضمانات الثلاث}

\en{Stripe} بيكلّمنا رجوعًا عن طريق \en{webhooks}: رسائل \en{HTTP} بيبعتها لخادمنا
لمّا يحصل حدث (اكتمل تفويض، حصل استرداد، اكتمل صرف...). المشكلة: أي حد يقدر يبعت
طلب \en{HTTP} مزيّف يدّعي إنّه من \en{Stripe}. عشان كده لازم \textbf{نتحقّق من
التوقيع} قبل ما نصدّق أي رسالة. ده بيحصل في \code{ParseWebhook}:

\begin{english}
\begin{lstlisting}[language={[Sharp]C}]
var stripeEvent = EventUtility.ConstructEvent(
    payload,
    signatureHeader,
    webhookSecret,
    throwOnApiVersionMismatch: false);
// throws StripeException if the signature is invalid -> we reject the request

switch (stripeEvent.Data?.Object)
{
    case PaymentIntent pi: paymentIntentId = pi.Id; amountCents = pi.Amount; break;
    case Charge ch:        chargeId = ch.Id; paymentIntentId = ch.PaymentIntentId; break;
    case Payout po:        payoutId = po.Id; break;
    // Account / Dispute cases ...
}
\end{lstlisting}
\end{english}

الدالة \code{EventUtility.ConstructEvent} بتاخد الحمولة الخام والتوقيع والسرّ
المشترك (\code{webhookSecret})، وتتأكّد إنّ التوقيع مطابق. لو مش مطابق، بترمي
استثناء ونرفض الطلب فورًا. بعد التحقّق، بنستخرج الكائن المناسب (نيّة دفع، فاتورة،
صرف...) عشان المُعالِج يقدر يطابقه مع صفّ الـ \code{Payment} الصحيح.

وعشان \en{Stripe} بيعيد إرسال نفس الحدث لحدّ 3 أيّام لو ما ردّيناش بنجاح، لازم
\textbf{نتجنّب معالجته مرّتين}. ده بيحصل بإزالة التكرار عبر
\code{StripeWebhookEvent.StripeEventId}:

\begin{english}
\begin{lstlisting}[language={[Sharp]C}]
// Idempotency: Stripe retries the same event for up to 3 days.
if (await db.StripeWebhookEvents.AnyAsync(e => e.StripeEventId == request.EventId, ct))
{
    logger.LogInformation("Webhook {EventId} already processed, skipping.", request.EventId);
    return true;
}
\end{lstlisting}
\end{english}

قبل ما نعمل أي حاجة، بنسأل: «الحدث ده اتعالج قبل كده؟». لو أيوة، بنتجاهله
ونرجّع نجاح. حقل \code{StripeEventId} عليه فهرس فريد في الداتابيز، فحتى لو
اتبعت نفس الحدث مرّتين في نفس اللحظة، واحد بس بينجح والتاني بيتعامل معاه كنجاح
مكرّر.

\begin{secnote}[زاوية الأمان: ثلاث ضمانات تحمي مسار الفلوس]
\begin{itemize}
\item \textbf{مفاتيح عدم التكرار} (\en{idempotency keys}): ما فيش خصم مزدوج ولا
  صرف مزدوج مهما اتكرّرت المحاولات.
\item \textbf{القبض اليدويّ} (\en{manual capture}): الفلوس ما بتتاخدش إلا لمّا
  الخدمة تتأكّد فعلًا (المستشار قبل).
\item \textbf{\en{webhooks} موقَّعة + إزالة تكرار الأحداث}: ما فيش حدث دفع مزيّف
  أو معاد إرساله يقدر يخدعنا.
\end{itemize}
وفوق ده كلّه: \textbf{بيانات الكروت لا تلمس خوادمنا أبدًا}. المتصفّح بيكلّم
\en{Stripe} مباشرة عن طريق الـ \en{client secret}، وإحنا بس بنشوف معرّف نيّة
الدفع (\en{intent id}). فحتى لو تسرّبت قاعدة بياناتنا، ما فيش أرقام كروت تتسرّب.
\end{secnote}

\begin{qa}
\textbf{سؤال متوقّع: «فين الفلوس ما بين الحجز والجلسة؟»} \\
الإجابة: مُفوَّضة ومحجوزة على بطاقة الطالب عبر نيّة دفع بنمط \en{manual capture}.
بنقبضها بس لمّا المستشار يقبل؛ ولو رفض أو انتهت المهلة بنفكّ الحجز؛ ولو حصل
\en{no-show} أو إلغاء مؤهَّل، \code{RefundCalculatorService} بيقرّر الاسترداد.
وصافي الأرباح بيوصل للمستشار عبر صرف \en{Stripe Connect}.
\end{qa}

\subsection{المهامّ الخلفيّة (Hangfire)}

كتير من الشغل في النظام مش بيحصل لحظة طلب المستخدم، لكن بيحصل «لوحده» في الخلفية
على فترات منتظمة. إحنا بنستخدم \en{Hangfire} لجدولة المهامّ دي. كل المهامّ
متحكَّم فيها بمفتاح إعداد واحد: \code{Hangfire:Enabled}؛ لو اتقفل، ما بيتشغّلش أي
جدول (مفيد في بيئات التطوير والاختبار).

الجدول ده بيلخّص المهامّ المُجدوَلة الفعليّة في النظام (13 مهمّة):

\begin{center}
\begin{tabular}{L{4.9cm} L{2.5cm} L{5.1cm}}
\toprule
\textbf{\color{spPrimary}المهمّة} & \textbf{\color{spPrimary}الجدولة} & \textbf{\color{spPrimary}الغرض} \\
\midrule
\code{StripePayoutJob} & يوميًّا 02:00 & صرف صافي الأرباح للمستشارين دفعة واحدة \\
\code{MeetingNoShowSweepJob} & كل 15 دقيقة & رصد غياب أحد الطرفين وفتح مسار الاسترداد \\
\code{SessionExpiryJob} & كل 15 دقيقة & إنهاء الحجوزات اللي فاتت مهلة ردّ المستشار \\
\code{CompletionJob} & كل 15 دقيقة & إغلاق الجلسات المكتملة تلقائيًّا \\
\code{BookingReminderJob} & كل 15 دقيقة & تذكير الطالب والمستشار قبل الجلسة \\
\code{DeadlineReminderJob} & يوميًّا 09:00 & تنبيه الطلاب بقرب انتهاء مواعيد المنح \\
\code{NotificationDispatcherJob} & يوميًّا 10:00 & تذكير بالمسودّات غير المُرسَلة \\
\code{ScholarshipAutoCloseJob} & كل ساعة & إغلاق المنح اللي فات موعدها تلقائيًّا \\
\code{ReviewTimeoutRefundJob} & يوميًّا 05:00 & استرداد رسوم مراجعة تأخّرت عن مهلتها \\
\code{DataExportJob} & كل ساعة & تجهيز تصدير بيانات المستخدم (\en{GDPR}) \\
\code{DataDeleteJob} & يوميًّا 03:00 & حذف/تعمية بيانات المستخدم (\en{GDPR}) \\
\code{RedactionAuditSamplingJob} & شهريًّا & عيّنات تنقية للتحقّق من حجب الـ \en{PII} \\
\code{IntegrityCheckJob} & يوميًّا 04:00 & فحوص سلامة للصفوف اليتيمة والعالقة \\
\bottomrule
\end{tabular}
\end{center}

نلاحظ إنّ المهامّ دي مش بس تذكيرات؛ فيها منها اللي بيمسّ الفلوس مباشرة
(\code{StripePayoutJob}، واسترداد المراجعة)، ومنها اللي بيخدم الالتزام القانونيّ
(\en{GDPR})، ومنها اللي بيراقب صحّة النظام نفسه (\code{IntegrityCheckJob} بيعدّ
الدفعات اليتيمة والـ \en{webhooks} العالقة). كلها مبنيّة عشان تكون
\textbf{آمنة للتكرار}: لو المهمّة اتشغّلت مرّتين، مفيش ضرر — نفس مبدأ عدم التكرار
اللي شفناه في مسار الدفع.

\begin{why}
ليه نفصل الشغل ده في الخلفية بدل ما نعمله لحظة الطلب؟ عشان تجربة المستخدم تفضل
سريعة (ما نخلّيهوش يستنّى صرف فلوس أو تصدير بيانات)، وعشان الشغل الدوريّ (زيّ
إغلاق المنح المنتهية) ما ليهوش «طلب مستخدم» أصلًا يشغّله — لازم حد يفتكره ويشغّله،
والـ \en{Hangfire} هو ده «الحد».
\end{why}

\begin{summary}[الخلاصة في سطور]
\begin{itemize}
\item النظام بيستخدم \en{authorize-then-capture} (شبه \en{escrow}) على \en{Stripe}:
  فلوس الطالب \textbf{محجوزة لا مأخوذة} لحدّ ما المستشار يقبل.
\item المنفذ \code{IStripeService} له محوّلان: \code{StripeService} الحقيقي و
  \code{StubStripeService} للتطوير — نبدّل من غير ما نلمس منطق الحجز.
\item دورة الحياة: حجز + تفويض (\en{manual capture}، سنتات صحيحة، مفتاح عدم تكرار)
  ثمّ قبول = قبض ثمّ رفض/انتهاء = إفراج ثمّ \en{no-show}/إلغاء = استرداد بحساب
  \code{RefundCalculatorService} ثمّ صرف عبر \code{StripePayoutJob}.
\item تقسيم الأرباح بالسنتات: \code{AmountCents} تساوي \code{ProfitShareAmountCents}
  (المنصّة) + \code{PayeeAmountCents} (المستشار)، من قاعدة ماليّة قابلة للتهيئة،
  بفهرس فريد مُصفّى يضمن قاعدة فعّالة واحدة لكل نوع.
\item الـ \en{webhook} بيتحقّق من توقيع \en{Stripe} ويزيل التكرار عبر
  \code{StripeEventId}. ثلاث ضمانات: عدم التكرار + القبض اليدويّ + \en{webhooks}
  موقَّعة. وبيانات الكروت \textbf{ما بتلمسش خوادمنا}.
\item 13 مهمّة \en{Hangfire} خلفيّة (تذكيرات، صرف، \en{no-show}، إغلاق منح،
  \en{GDPR}، تنقية، فحوص سلامة) متحكَّم فيها بـ \code{Hangfire:Enabled}.
\end{itemize}
\end{summary}
